% packages.tex

% Sprache 'english', 'ngerman', ...
\usepackage[ngerman]{babel}

% Schriftart
\usepackage{fontspec}
\setmainfont{Times New Roman}

% Bessere Absatzformatierung
\usepackage{microtype}
\sloppy
\hyphenpenalty=500
\exhyphenpenalty=50

% Keine Überschriften am Ende der Seite
\usepackage{needspace}
\let\oldsection\section
\renewcommand{\section}{\needspace{5\baselineskip}\oldsection}

% Zitate
\usepackage[
backend=bibtex,
style=authoryear	%hier: Chicago
]{biblatex}
\addbibresource{/home/marcel/references.bib}	%.bib datei
% Für Kurzzitate, wie z.B. Wittgenstein oder Aristoteles
\DeclareCiteCommand{\shortcite}[\mkbibparens]
{\usebibmacro{prenote}}
{\printfield{shorttitle}}
{\multicitedelim}
{\addspace\printfield{postnote}}
\DeclareFieldFormat{postnote}{#1} %entfernt p. vor Seitenzahl

% Seitenlayout
\usepackage[a4paper, left=2.54cm, right=2.54cm, top=2.54cm, bottom=2.54cm]{geometry}

% 1,5 Zeilen Abstand
\usepackage{setspace}
\setstretch{1.5}

\setlength{\parindent}{0pt} % Kein Einzug am Anfang des Absatzes
\setlength{\parskip}{6pt} % Abstand zwischen Absätzen

% Formatierung der Kopf- und Fußzeile
\usepackage{fancyhdr}
\pagestyle{fancy}
\fancyhf{}
\fancyfoot[C]{\thepage}
\renewcommand{\headrulewidth}{0pt}
\renewcommand{\footrulewidth}{0pt}

% Styling der Überschriften
\usepackage{titlesec}
\titleformat{\section}
{\normalfont\Large\bfseries}{\thesection.}{1em}{}
\titleformat{\subsection}
{\normalfont\large\bfseries}{\thesubsection.}{1em}{}
\titleformat{\subsubsection}
{\normalfont\normalsize\bfseries}{\thesubsubsection.}{1em}{}
\titlespacing*{\section}{0pt}{1.5\baselineskip}{0.5\baselineskip}
\titlespacing*{\subsection}{0pt}{1\baselineskip}{0.25\baselineskip}

% Inhaltsverzeichnis
\usepackage{tocloft}
% Abstand unter "Inhaltsverzeichnis" wie bei normalen Sections
\renewcommand{\cftaftertoctitle}{\par\vspace{0.5\baselineskip}}

% Schrift wie im Dokument
\renewcommand{\cftsecfont}{\normalfont}
\renewcommand{\cftsubsecfont}{\normalfont}
\renewcommand{\cftsubsubsecfont}{\normalfont}

% Seitenzahlen nicht fett
\renewcommand{\cftsecpagefont}{\normalfont}
\renewcommand{\cftsubsecpagefont}{\normalfont}
\renewcommand{\cftsubsubsecpagefont}{\normalfont}

% Punktlinien zwischen Titel und Seitenzahl (klassischer Word-Stil)
\renewcommand{\cftdotsep}{1}
% Einheitliches Punktlinien-Verhalten für alle Ebenen
\renewcommand{\cftsecleader}{\cftdotfill{\cftdotsep}}
\renewcommand{\cftsubsecleader}{\cftdotfill{\cftdotsep}}
\renewcommand{\cftsubsubsecleader}{\cftdotfill{\cftdotsep}}

% Abstände zwischen Einträgen
\setlength{\cftbeforesecskip}{2pt}
\setlength{\cftbeforesubsecskip}{2pt}
\setlength{\cftbeforesubsubsecskip}{2pt}

% Einrückungen (Hierarchie)
\setlength{\cftsecindent}{0pt}
\setlength{\cftsubsecindent}{1.5em}
\setlength{\cftsubsubsecindent}{3em}

% Nummernbreiten
\setlength{\cftsecnumwidth}{1.8em}
\setlength{\cftsubsecnumwidth}{2.6em}
\setlength{\cftsubsubsecnumwidth}{3.4em}

% Punkt nach der Nummer im Inhaltsverzeichnis
\renewcommand{\cftsecaftersnum}{.}
\renewcommand{\cftsubsecaftersnum}{.}
\renewcommand{\cftsubsubsecaftersnum}{.}

% Überschrift "Inhaltsverzeichnis" wie eine Section
\renewcommand{\cfttoctitlefont}{\Large\bfseries}
\renewcommand{\cftaftertoctitle}{\par\vspace{1\baselineskip}}

% Für eingerückte Blockzitate
\usepackage{csquotes}
\renewenvironment{displayquote}
{\list{}{\leftmargin=1.5cm \rightmargin=1.5cm}\item\relax}
{\endlist}
